\section{Inverse Reinforcement Learning}

\begin{frame}{}
    \LARGE The Reward Problem: \textbf{Inverse RL}
\end{frame}

\begin{frame}{Inverse RL: Given Behavior, Learn the Reward}
\begin{itemize}
    \setlength{\itemsep}{0.6em}
    \item \textbf{Forward RL} (everything so far): given a reward $R(s,a)$, learn a policy $\pi^*$ that maximizes it
    \item \textbf{Inverse RL} (today): given (expert) behavior $\tau^E \sim \pi^E$, learn a reward $R_\theta$ that \emph{explains} it
    \[
        \pi^E \;\;\Longrightarrow\;\; R_\theta
    \]
\end{itemize}
\end{frame}

\begin{frame}{Why Inverse RL?}
\begin{itemize}
    \setlength{\itemsep}{0.65em}
    \item In many domains, rewards are far easier to \emph{demonstrate} than to specify in closed form
    \item Driving: hard to write ``polite, safe, fast'' as a scalar; easy to show
    \item Manipulation, dialogue, locomotion: same story
    \item \textcolor{red}{\textbf{This is the strongest form of the reward problem}} --- the reward function itself is unknown.
\end{itemize}
\end{frame}

\begin{frame}{The Ill-Posedness of IRL}
\begin{itemize}
    \setlength{\itemsep}{0.6em}
    \item \emph{Many} reward functions rationalize the same expert behavior
    \begin{itemize}
        \item Trivially: $R(s,a) \equiv 0$ makes every policy optimal
        \item Potential-based shaping: $R$ and $R + \gamma\Phi(s') - \Phi(s)$ are policy-equivalent
    \end{itemize}
    \item We need a \emph{principle} to pick one reward out of the equivalence class
\end{itemize}
\end{frame}

\begin{frame}{Two Structural Fixes for the Ambiguity}
\begin{itemize}
    \setlength{\itemsep}{0.75em}
    \item \textbf{(a) Max-margin IRL} \\
          \emph{Ng \& Russell, ICML 2000;\ Abbeel \& Ng, ICML 2004} \\
          The expert's reward should beat alternatives by a wide margin
    \item \textbf{(b) Maximum-entropy IRL} \\
          \emph{Ziebart et al., AAAI 2008} \\
          Pick the reward whose induced policy distribution has \emph{maximum entropy}, subject to matching the expert's feature counts --- a unique, well-posed choice
    \item We focus on (b) --- standard fix today, and generalizes to modern preference learning
\end{itemize}
\end{frame}

\begin{frame}{Maximum-Entropy IRL: The Principle}
\textbf{Among all trajectory distributions that match the expert's behaviour, pick the one with maximum entropy.}
\vspace{0.3em}
\begin{itemize}
    \setlength{\itemsep}{0.4em}
    \item Written out, that is a \emph{constrained} problem:
    \[
        \max_{P} \; H(P)
        \quad \text{s.t.} \quad
        \mathbb{E}_{P}\big[\text{feature counts}\big] = \text{expert's feature counts}
    \]
    \item \textcolor{red}{Solving it gives the exponential form} --- higher-reward trajectories become exponentially more likely:
    \[
        P_\theta(\tau) \;\propto\; \exp\!\Big(\textstyle\sum_t R_\theta(s_t, a_t)\Big)
    \]
    \item \textcolor{gray}{So there is no entropy \emph{term} to point at in that formula --- the entropy has already been solved out, and the $\theta$ are the Lagrange multipliers of the matching constraint}
    \item Train $\theta$ by maximum likelihood on expert trajectories: $\;\max_\theta \sum_{\tau^E} \log P_\theta(\tau^E)$
    \item Exactly \emph{one} distribution satisfies both conditions --- which is what makes the choice unique
\end{itemize}
\end{frame}

\begin{frame}{MaxEnt IRL: The Day-6 Connection}
\vspace{0.4em}
\begin{center}
\textcolor{red}{\textbf{The same max-entropy principle that justified SAC's $-\log\pi$ bonus.}}
\end{center}
\vspace{0.8em}
\begin{itemize}
    \setlength{\itemsep}{0.7em}
    \item In SAC (Day 6): adds exploration to the policy
    \item In MaxEnt IRL: resolves reward ambiguity
    \item \textbf{Same principle, two different jobs.}
    \item \textcolor{gray}{Different mechanics, though: SAC \emph{adds} $\alpha\mathcal{H}(\pi)$ to the objective, so you can point at the term. MaxEnt IRL imposes entropy as the \emph{selection rule}, which is why it disappears into the exponential form}
    \item Cite: \textbf{Ziebart, Maas, Bagnell \& Dey, AAAI 2008} --- \emph{Maximum Entropy Inverse Reinforcement Learning}
\end{itemize}
\end{frame}

% Bridges the 2008 -> today gap: how MaxEnt IRL became deep-learning-scale.
\begin{frame}{From 2008 to Today: How IRL Scaled}
\textbf{The catch with MaxEnt IRL as stated:} it needs hand-designed feature counts and a tractable normaliser --- fine for a tabular gridworld, hopeless for pixels.
\vspace{0.3em}
\begin{itemize}
    \setlength{\itemsep}{0.4em}
    \item \textbf{Guided Cost Learning} (Finn et al., 2016): estimate that normaliser by \emph{sampling} instead of enumerating $\Rightarrow$ neural-network rewards
    \item \textcolor{red}{Then the surprise:} running a \textbf{GAN} on this problem optimises \emph{the same objective} as MaxEnt IRL --- the discriminator plays the role of the reward, the generator the policy
    \item \textbf{GAIL} (Ho \& Ermon, 2016) and \textbf{AIRL} (Fu et al., ICLR 2018) are that idea made practical: learn reward and policy \emph{adversarially}, straight from demonstrations
    \item \textcolor{gray}{Meanwhile robotics largely went the other way --- large-scale \emph{behaviour cloning} (Diffusion Policy, VLAs) is now the dominant way to learn from demonstrations there}
    \item \textcolor{red}{\textbf{IRL's centre of gravity moved to language models}} --- which is where the next two slides are heading
\end{itemize}
\end{frame}

\begin{frame}{Preference-Based Reward Learning: The Setup}
\textbf{What if we have neither demonstrations nor a written reward, only human \emph{comparisons}?}
\vspace{0.6em}
\begin{itemize}
    \setlength{\itemsep}{0.55em}
    \item Show two trajectories $\tau^A, \tau^B$; ask a human which is better
    \item Model the preference with a \textbf{Bradley--Terry} likelihood under a learned reward $R_\theta$:
    \[
        P(\tau^A \succ \tau^B) \;=\;
        \frac{\exp\!\big(\sum_t R_\theta(s_t^A, a_t^A)\big)}
             {\exp\!\big(\sum_t R_\theta(s_t^A, a_t^A)\big)
            + \exp\!\big(\sum_t R_\theta(s_t^B, a_t^B)\big)}
    \]
\end{itemize}
\end{frame}

\begin{frame}{Preference-Based Reward Learning: The Pipeline}
\begin{enumerate}
    \setlength{\itemsep}{0.6em}
    \item Collect a dataset of human-labeled preferences over trajectory pairs
    \item Train $R_\theta$ by cross-entropy on the Bradley--Terry likelihood
    \item Run any standard RL algorithm (e.g.\ PPO from Day 4) against $R_\theta$
\end{enumerate}
\vspace{0.7em}
\begin{itemize}
    \item Cite: \textbf{Christiano, Leike, Brown, Martic, Legg \& Amodei, NeurIPS 2017}
\end{itemize}
\end{frame}

\begin{frame}{The Day-9 Bridge}
\vspace{0.6em}
\begin{center}
\textcolor{red}{\textbf{This exact pipeline + a language-model policy $=$ RLHF.}}
\end{center}
\vspace{0.9em}
\begin{itemize}
    \setlength{\itemsep}{0.65em}
    \item The reward model in RLHF is the $R_\theta$ from the previous slide
    \item The only thing that changes is the policy class: $\pi_\theta$ becomes a large language model
    \item The optimizer (PPO) and the pipeline (SFT $\to$ RM $\to$ PPO) are the same machinery
    \item \textcolor{red}{Read the other way:} training a reward model \emph{is} an inverse-RL problem --- inferring a reward from human data rather than writing one down
    \item \textcolor{gray}{And the loop closes: 2025 work runs adversarial IRL on expert chain-of-thought to recover \emph{process-level} reasoning rewards, instead of imitating the traces}
    \item Day 9 picks up exactly here.
\end{itemize}
\end{frame}

\begin{frame}{The Three Faces, Resolved}
\begin{center}
\renewcommand{\arraystretch}{1.6}
\small
\begin{tabular}{p{0.42\textwidth}p{0.50\textwidth}}
\hline
\textbf{If you can\ldots} & \textbf{Then use\ldots} \\
\hline
\ldots write the reward & \emph{Forward RL} (Days 1--6) \\
\ldots demonstrate the reward & \emph{Inverse RL} (MaxEnt IRL) \\
\ldots compare outcomes & \emph{Preference-based reward learning} \\
                       & (Day 9: RLHF) \\
\hline
\end{tabular}
\end{center}
\vspace{1em}
\begin{itemize}
    \item One unifying view: every method on the board is a different way to recover the optimal policy when the reward is unknown, partially known, or hard to specify
\end{itemize}
\end{frame}
